% !TeX root = ../main.tex
%%% ---------------------------------------------------------------------------
%%% 国内外研究现状与存在问题
%%%
%%% 目标页数：2~3 页。
%%% 这一节的成败标准是：委员会看完能不能说出「所以现有方法差在哪」。
%%% 常见失误是把它写成文献流水账——列了 20 篇，但没有一句话说清共同的局限。
%%%
%%% 一个有效的做法：用一张表把现有路线按「解决了什么 / 没解决什么」摆开，
%%% 最后一行是自己的工作。委员会一眼就看到 gap 在哪。
%%% ---------------------------------------------------------------------------
\section{研究现状与存在问题}

\begin{frame}{现有路线各自解决了一半问题，但都没解决特征错位}
  \begin{table}
    \centering\small
    \begin{tabular}{lccc}
      \toprule
      {技术路线} & {多尺度表示} & {空间自适应} & {层级对齐} \\
      \midrule
      图像金字塔        & $\checkmark$ & --         & --         \\
      特征金字塔 FPN    & $\checkmark$ & --         & --         \\
      注意力增强骨干    & $\checkmark$ & $\checkmark$ & --         \\
      \midrule
      \textbf{本文}     & $\checkmark$ & $\checkmark$ & $\checkmark$ \\
      \bottomrule
    \end{tabular}
  \end{table}
  \takeaway{下采样引入的感受野偏移一直被当作可忽略量，小目标上它不可忽略}

  \note{这张表是防守用的。委员会若问「和某某工作有什么区别」，
        直接指第三列或第四列。备用页有逐篇对比。}
\end{frame}

\begin{frame}{存在的两个具体问题}
  \begin{enumerate}
    \setlength{\itemsep}{1.2em}
    \item<1-> \textbf{融合权重在空间上均匀}\\
      {\small 背景区域与小目标区域被同等对待，浅层的高分辨率信息
       在前者上被浪费、在后者上又不够。}
    \item<2-> \textbf{相邻层级的感受野未对齐}\\
      {\small 下采样引入亚像素偏移，逐元素相加造成特征错位，
       误差在小目标上被成倍放大。}
  \end{enumerate}
  \vspace{0.8em}
  \onslide<3->{\takeaway{本文的两块工作分别针对这两个问题}}
\end{frame}
